\section{Agents for Science：把``AI 做科研''变成可测量现象}

第三部分最重要的贡献是方法学：不是再展示单个 demo，而是组织一个真实会议，把投稿、评审、人机分工数据公开，进而把大量讨论变成可分析数据。

Zou 把会议规则设为``AI 必须是一作并完成主要工作，人类可协作''。这在当下主流会议制度下很激进，但正因激进，才有机会系统观察 AI-first 研究范式的优劣。

\subsection{会议设计与样本规模}

课程报告中，会议收到 300+ 投稿，录用 48 篇，覆盖工程、物理、医学健康、社会科学等方向。每篇都要提交 AI involvement checklist，按 hypothesis、experimental design、data analysis、writing 四阶段标注人机贡献比。

\begin{importantbox}{为何这套会议设计有研究价值}
\begin{itemize}
    \item 规则显式化：避免``用了 AI 但不披露''导致的数据偏差。
    \item 过程结构化：将贡献按科研阶段拆解，而不是笼统自报。
    \item 结果可对照：AI 评审结果可与人类 spot-check 交叉比较。
\end{itemize}
\end{importantbox}

图\ref{fig:conference-questions} 和图\ref{fig:conference-reviews} 对应该部分关键帧。\footnote{来源：Berkeley RDI 课程视频，关键帧时间戳 00:44:40。}\footnote{来源：Berkeley RDI 课程视频，关键帧时间戳 00:52:30。}

\begin{figure}[H]
\centering
\includegraphics[width=0.9\textwidth]{frame-04-benchmark.jpg}
\caption{Agents for Science 提出的核心问题}
\label{fig:conference-questions}
\end{figure}

\begin{figure}[H]
\centering
\includegraphics[width=0.9\textwidth]{frame-05-conference.jpg}
\caption{LLM 评审示例：可发现错误，也会出现谄媚式评价}
\label{fig:conference-reviews}
\end{figure}

\begin{knowledgebox}{读图：问题清单与评审样例组成一套观测设计}
第一张图定义要测量的现象：原创性、人机分工和 AI 评审质量；第二张图给出评审能力的两面性，既能抓到统计不一致，也会生成谄媚式高评价。把二者合起来看，会议不是证明 AI 已能独立科研，而是为协作行为、错误和偏置提供可分析样本。
\end{knowledgebox}

\subsection{人机协作模式观察}

课程结论之一是：录用论文整体上保留了更多人类介入，尤其在假设提出和实验设计前段；而在后段（数据分析、写作）更容易出现 AI 主导。这与实际科研风险分布一致，前段决策错误会造成后续系统性偏航，因此人类更倾向把控上游环节。

\begin{knowledgebox}{四阶段协作启示}
\begin{itemize}
    \item Hypothesis：人类主导方向选择，AI 提供候选与反例。
    \item Experiment Design：人机共设约束，避免不可执行方案。
    \item Data Analysis：AI 负责高吞吐执行，人类做统计审查。
    \item Writing：AI 可承担初稿整合，人类负责论证强度与责任边界。
\end{itemize}
\end{knowledgebox}

\subsection{AI 评审的异质性：保守、激进与居中}

课程中使用 GPT-5、Gemini 2.5 Pro、Claude Sonnet 4 作为评审 Agent。观察到明显评分风格差异：GPT-5 相对保守，Gemini 相对更积极，Claude 在展示数据里更接近人类评分分布。这说明``AI 评审''不是单一实体，而是带模型风格偏置的评估系统。

\begin{warningbox}{评审自动化的风险不在于``会不会打分''，而在于``评分偏置可否校正''}
如果会议把单模型评审结果直接当最终判断，可能把某一类风格偏好固化成制度偏见。更稳妥的是多模型委员会 + 人类抽检 + 评分校准曲线共同作用。
\end{warningbox}

\subsection{引用幻觉检测：可量化的真实痛点}

会议还上线了自动引用核验流程：抽取每篇参考文献标题并检索匹配。报告显示约 44\% 投稿参考文献可完全核验，约 56\% 至少出现一条不可核验引用（多数是少量错误，也有极端样本）。

这组数字的重要性在于：它把``AI 写作可能胡编引用''从坊间印象变成结构化证据，也为后续治理提供可操作指标。

\begin{importantbox}{引用核验应成为 AI-first 论文的默认基础设施}
\begin{itemize}
    \item 投稿阶段自动核验并返回 warning 列表。
    \item 评审阶段强制作者对不可核验引用逐条回应。
    \item 录用后公开核验报告，形成可追责记录。
\end{itemize}
\end{importantbox}

\subsection{本章小结}

Agents for Science 的价值在于把争论变成数据：我们不仅看到 AI 可以做什么，也看到它在哪些环节必须受约束、哪些环节需要人类补位。

